\documentclass[aps,preprint]{revtex4}%
\usepackage{amsfonts}
\usepackage{amsmath}
\usepackage{amssymb}
\usepackage{graphicx}%
\setcounter{MaxMatrixCols}{30}
%TCIDATA{OutputFilter=latex2.dll}
%TCIDATA{Version=4.10.0.2363}
%TCIDATA{CSTFile=revtex4.cst}
%TCIDATA{Created=Friday, May 16, 2003 15:55:00}
%TCIDATA{LastRevised=Sunday, January 25, 2004 08:59:15}
%TCIDATA{<META NAME="GraphicsSave" CONTENT="32">}
%TCIDATA{<META NAME="DocumentShell" CONTENT="Articles\SW\REVTeX 4">}
%TCIDATA{Language=American English}
\newtheorem{theorem}{Theorem}
\newtheorem{acknowledgement}[theorem]{Acknowledgement}
\newtheorem{algorithm}[theorem]{Algorithm}
\newtheorem{axiom}[theorem]{Axiom}
\newtheorem{claim}[theorem]{Claim}
\newtheorem{conclusion}[theorem]{Conclusion}
\newtheorem{condition}[theorem]{Condition}
\newtheorem{conjecture}[theorem]{Conjecture}
\newtheorem{corollary}[theorem]{Corollary}
\newtheorem{criterion}[theorem]{Criterion}
\newtheorem{definition}[theorem]{Definition}
\newtheorem{example}[theorem]{Example}
\newtheorem{exercise}[theorem]{Exercise}
\newtheorem{lemma}[theorem]{Lemma}
\newtheorem{notation}[theorem]{Notation}
\newtheorem{problem}[theorem]{Problem}
\newtheorem{proposition}[theorem]{Proposition}
\newtheorem{remark}[theorem]{Remark}
\newtheorem{solution}[theorem]{Solution}
\newtheorem{summary}[theorem]{Summary}
\newenvironment{proof}[1][Proof]{\noindent\textbf{#1.} }{\ \rule{0.5em}{0.5em}}
\begin{document}
\preprint{ }
\title[AGP\ Coh. States ]{Antisymmetrized Geminal Power Coherent States}
\author{Brian Weiner}
\affiliation{Department of Physics, Pennsylvania State University, DuBois PA 15801}
\keywords{}

\begin{abstract}


\end{abstract}
\date{01/19/2004}
\startpage{1}
\maketitle
\tableofcontents


\section{Introduction\label{intro}}

Isolated molecular systems do not interact with their environment and are thus
unobservable. In order to have a physical effect and be observed they must
respond to external probes, such as electromagnetic fields, and form a
composite interacting system whose evolution is described by a joint
Hamiltonian. The effect of radiation on the molecular subsystem can be
monitored by changes in the molecules density and reduced density matrices. If
the radiation is weak the behavior of the interaction is well described by the
linear response of the molecular system to the radiation, i.e. is linearly
proportional to the radiation strength, while if the radiation is intense the
molecular system responds in a nonlinear way. The interaction of the
electromagnetic field with molecules can be approximated in a standard semi
classical fashion by averaging over the photonic degrees of freedom to produce
effective one electron observables. This leads to a reduced Hamiltonian for
the molecular subsystem that includes perturbing effective time dependent one
body potentials which describe the interaction with the radiation field. In
the case of weak intensity radiation, where the linear response regime is
sufficient to describe the changes in the molecular system, one needs to only
consider responses that are first order in these observables.

The temporal behavior of a state of physical system is completely described by
the evolution of its full density matrix, which in term determines the
evolution of all of its reduced density matrices. These variations are
conveniently expressed in terms of various propagators, which thus provide a
comprehensive formulation of the response of the molecular subsystem to
external probes. As is well known exact solutions for quantum evolutions have
only been determined for simple model systems, and are not known for
multi-electron molecular systems that involve coulombic interactions between
electrons. Thus for practical applications the goal is to develop systematic,
cost effective, theoretically well founded and non arbitrary approximation
schemes, that are accurate enough to explain experimental results. A major
focus of the illustrious career of professor Osvaldo Goscinski has been the
analysis and construction of approximate propagators, that fulfill these
objectives. In particular the one-electron and polarization propagators, where
he effectively used the properties of super operators acting on operator
manifolds associated with a given reference state to generate various orders
of approximation \cite{Goscinski73}\cite{Pickup73}\cite{Goscinski69}%
\cite{Goscinski70}\cite{Weiner80a}\cite{Goscinski80}\cite{Howat77a}. In this
article I will concentrate, for illustrative purposes, on the polarization
propagator for a molecular system with a fixed even number of electrons at a
fixed geometry. This propagator describes the linear response of the First
Order Reduced Density Matrix (FORDM) to electromagnetic radiation.

The polarization propagator is a hermitian matrix of time dependent response
functions of the matrix elements of the FORDM. If one takes the Laplace
transform of this propagator and expresses the resulting matrix in diagonal
form, in terms of normal mode operators, one can clearly observe that the
poles of this propagator are excitation energies and the associated residues
the transition probabilities for excitations generated by one particle
operators between stationary states, $\left\{  \left\vert \Psi_{k}%
\right\rangle \right\}  ,$ of the molecular system. If the initial molecular
reference state is stationary (usually the ground state, but not necessarily
so) then the poles are the excitation energies and the residues the transition
probabilities from that state. If the transition probability is 100\% for a
given one particle operator then this operator produces an excited state from
the initial stationary state and is called a state excitation operator. If
such an operator produces the excited state $\left\vert \Psi_{k}\right\rangle
,$ then its effect on the reference state $\left\vert \Psi_{0}\right\rangle $
is the same as the $N$-particle bra-ket operator $\left\vert \Psi
_{k}\right\rangle \left\langle \Psi_{0}\right\vert .$

Approximations to the polarization propagator are often based on following criteria,

\begin{itemize}
\item all one particle operators that diagonalize the propagator matrix are
state excitation operators,

\item their adjoints annihilate the reference state,

\item the reference state is stationary with respect to at least one particle variations.
\end{itemize}

If an approximation satisfies these conditions one says that it is a
consistent approximation, however, in general, usually one or more of these
conditions are not satisfied and the approximation is inconsistent. Some
approximations can even lead to propagators that do not correspond to a valid
$N$-particle reference state and thus are not even $N$-representable.

An important point of note is that it is not necessary that the normal mode
one-particle operators actually produce excited states from the reference
state, in fact, in the exact case they generally do not. Hence there is thus
some level of ambiguity about polarization propagator approximations, if one
is trying to obtain an approximation to a one-particle linear response
function then the annihilation condition is not necessary, though stationarity
and $N$-representability are desirable, while if one is asserting that the
polarization propagator is a one-particle approximation to the full
$N$-particle propagator and is used to produce excited states, as well as
excitation energies and transition probabilities, then the annihilation
condition is necessary for consistency.

In the late 1970's it was found \cite{Weiner80}\cite{Weiner80a}%
\cite{Goscinski80} \cite{Linderberg77}\cite{Ohrn79} that an Antisymmetrized
Geminal Power (AGP) \cite{Coleman63} State was the most general $2M$ electron
state that was annihilated by a maximal set of one electron operators, (for
odd number of electrons Generalized AGP (GAGP) states \cite{Coleman65} are
needed but their properties can be derived from AGP states). The $2M$-electron
AGP $\left\vert g^{M}\right\rangle $ state is an antisymmetrized product of
$M$ two particle geminals
\begin{equation}
\left\vert g\right\rangle =\sum_{1\leq j<k\leq s}g_{jk}\left\vert \varphi
_{j}\varphi_{k}\right\rangle
\end{equation}
where $\left\{  g_{ij}\right\}  $ are complex coefficients and $\left\{
\left.  \left\vert \varphi_{j}\right\rangle \right\vert 1\leq j\leq
r=2s\right\}  $ is an orthonormal basis for the one particle Hilbert space
$\mathcal{H}^{1}.$ However even when the reference state was an energy
optimized AGP state the polarization propagator approximation failed to be
consistent, in that the adjoints of the normal mode excitation operators did
not annihilate the reference state, although the numerical results obtained
from the application of this approximation scheme agreed well with
experimental values of excitation energies and dipole induced transition
probabilities, i.e. oscillator strengths, for example \cite{Sangfelt85}%
\cite{Sangfelt84}\cite{Elander83}\cite{Goscinski82}\cite{Weiner83}%
\cite{Sangfelt81} \cite{Weiner87}\cite{Sangfelt87}\cite{Ohrn85}\cite{Weiner85}%
\cite{Weiner84a}\cite{Jensen82}\cite{Jorgen83}.

It turns out that the annihilation condition can be naturally expressed in
terms of group Coherent States (CS) \cite{Perelomov_book}\cite{Klauder}%
\cite{Gilmore_book}, where the reference state corresponds to a lowest or
highest weight state associated with a representation of the group. The basic
ingredients needed for the construction of a family of group CS's in a vector
space is the existence of an irreducible representation of a group and a one
dimensional representation of a subgroup. The subgroup is called the isotropy
or invariance group of the CS's and the family of CS's is generated by the
action of the coset formed by factoring the group with respect to (wrt) this
subgroup on a reference state. In the study of electronic systems it is
natural to consider the real, compact, \cite{Encyclo} Lie group $U(r)$ of
unitary transformations of the one electron state space $\mathcal{H}^{1}$ (of
dimension $r=2s)$, that has irreducible representations on all $N$-electron
state spaces $\mathcal{H}^{N}$. The family of CS's in this case are nonlinear
manifolds of the $N$-electron Hilbert space, $\mathcal{H}^{N},$ generated by
cosets formed by one-particle operators.

One can construct approximate propagators based on stationary states of the
restriction of the Time Dependent Schr\"{o}dinger Equation (TDSE) to these
manifolds. Solutions to such restrictions of the TDSE are obtained by using
the Time Dependent Variational Principle (TDVP) \cite{Kramer81} or
equivalently the Stationary Phase Approximation (SPA) \cite{Schulman81} to the
path integral expressed in terms of paths on this CS manifold. The resulting
Equations of Motion (EOM)\ are those of a generalized \emph{classical }system
described by a curved phase space, a consequence of the fact that a family of
group coherent states forms a differentiable manifold \cite{Helgason78} whose
tangent spaces have an intrinsic symplectic structure \cite{Encyclo}.
Approximate stationary states of the molecular system then correspond to the
special stationary phase "paths" that are critical points of the corresponding
classical Hamiltonian. The linear response of these approximate stationary
states to external perturbing fields then involves the generators of the Lie
algebra of the group, and hence describes the response of the reduced density
matrix formed by these operators. Thus the response of the FORDM\ is
determined by generators of one-particle groups.

In this article we will describe how different coset decompositions of the one
particle unitary group $U(2s)$ correspond to different CS manifolds of
AGP\ states and how these manifolds can be parameterized by the coordinates of
the subspaces of the direct sum Lie algebra expansions adapted to these coset
decompositions. These CS\ manifolds form generalized curved phase spaces in
which classical EOM\ approximating the TDSE\ can be formulated, leading to a
classical linear response theory of time dependent perturbations that extends
and clarifies certain aspects of the quantum Random Phase Approximation and
propagator approximations. However due to space limitations these
considerations will be presented in future manuscripts.

In section \ref{liealg} we discuss the Lie algebra structure of Fermi-Fock
space, which is the direct sum of all of the state spaces $\left\{
\mathcal{H}^{N}\right\}  $, in section \ref{one particle groups} we will
describe the possible decompositions of the Lie algebra $u\left(  2s\right)  $
of $U(2s)$, the one particle group associated with $\mathcal{H}^{1},$ and thus
the factorizations of $U(2s)$ that lead to manifolds of coherent states in
$\mathcal{H}^{N}$ that can be based on $U(2s)$. In section \ref{cohstate} we
study in detail the form of AGP\ coherent states. We conclude with a
discussion in section \ref{discussion}.

\section{Lie Algebras\label{liealg}}

The most convenient representation of the generators of the Lie algebra of
$U(2s)$ for fermionic quantum mechanics is by second quantized operators
acting in Fermi-Fock space, $\mathcal{H}_{F}$, which is defined to be the
direct sum
\begin{equation}
\mathcal{H}_{F}=%
%TCIMACRO{\tbigoplus \limits_{0\leq N\leq r}}%
%BeginExpansion
{\textstyle\bigoplus\limits_{0\leq N\leq r}}
%EndExpansion
\mathcal{H}^{N}%
\end{equation}
where
\begin{equation}
\mathcal{H}^{N}=%
%TCIMACRO{\tbigwedge }%
%BeginExpansion
{\textstyle\bigwedge}
%EndExpansion
\mathcal{H}^{1}%
\end{equation}
is the $N$ -fold antisymmetric tensor product of the one particle state space
$\mathcal{H}^{1}$, a space generated by an orthonormal basis of one particle
functions of position and spin
\begin{equation}
\left\{  \left.  \varphi_{j}\right\vert 1\leq j\leq r\right\}  .
\end{equation}
The Banach space \cite{Encyclo} of bounded operators acting in $\mathcal{H}%
_{F}$ is denoted by $\mathcal{B}\left(  \mathcal{H}_{F}\right)  ,$ and is
spanned by a second quantized operator basis
\begin{equation}
\left\{  \left.  a_{j_{1}}^{\dagger}\cdots a_{j_{p}}^{\dagger}a_{k_{q}}\cdots
a_{k_{1}}\right\vert 1\leq j_{1}<\cdots j_{p}\leq r;1\leq k_{1}<\cdots
k_{p}\leq r;1\leq p,q\leq r\right\}
\end{equation}
where the second quantized operators $\left\{  a_{j}^{\dagger},a_{j}\right\}
$ act on the zero particle vacuum $\left\vert \phi\right\rangle $ in the
following fashion%
\begin{align}
a_{j}\left\vert \phi\right\rangle  &  =0\nonumber\\
a_{j}^{\dagger}\left\vert \phi\right\rangle  &  =\left\vert \varphi
_{j}\right\rangle
\end{align}
and satisfy the Canonical Anti-commutation Relationships (CAR)
\begin{equation}
\left[  a_{j}^{\dagger},a_{k}\right]  _{+}=\delta_{jk}I
\end{equation}
The operator space $\mathcal{B}\left(  \mathcal{H}_{F}\right)  $ contains
realizations of many classical Lie groups, algebras and universal enveloping
algebras \cite{Encyclo}. In particular the generators
\begin{equation}
\left\{  a_{j}^{\dagger},a_{k},a_{j}^{\dagger}a_{k}-\frac{1}{2}\delta
_{jk}I,a_{j}^{\dagger}a_{k}^{\dagger},a_{j}a_{k};1\leq j\leq k\leq r\right\}
\end{equation}
satisfy the commutation relationships of the lie algebra $so\left(
2r+1,\mathbb{R}\right)  ,$ (see \cite{Gilmore_book} for example for the
definitions of classical lie groups and algebras), and generate a
representation of the universal enveloping algebra $so\left(  2r+1,\mathbb{R}%
\right)  $ in $\mathcal{B}\left(  \mathcal{H}_{F}\right)  ,$ which is fact all
of $\mathcal{B}\left(  \mathcal{H}_{F}\right)  .$ The subset of these
generators
\begin{equation}
\left\{  a_{j}^{\dagger}a_{k}-\frac{1}{2}\delta_{jk}I,a_{j}^{\dagger}%
a_{k}^{\dagger},a_{j}a_{k};1\leq j\leq k\leq r\right\}
\end{equation}
satisfy the commutation relationships of the lie algebra of $so\left(
2r,\mathbb{R}\right)  $ and generate a representation of the universal
enveloping algebra of $so\left(  2r,\mathbb{R}\right)  $ in $\mathcal{B}%
\left(  \mathcal{H}_{F}\right)  .$ The subset of these generators
\begin{equation}
\left\{  a_{j}^{\dagger}a_{k}-\frac{1}{2}\delta_{jk}I;1\leq j\leq k\leq
r\right\}
\end{equation}
satisfy the commutation relationships of the lie algebra $u\left(  r\right)  $
and generate a representation of the universal enveloping algebra of $u\left(
r\right)  $ in $\mathcal{B}\left(  \mathcal{H}_{F}\right)  $. The
representation of $so\left(  2r+1,\mathbb{R}\right)  $ is irreducible, while
the representations of $so\left(  2r,\mathbb{R}\right)  $ and $u\left(
r\right)  $ are reducible. The irreducible representations of $so\left(
2r,\mathbb{R}\right)  $ are obtained by restricting to the subspace of even or
odd number states in $\mathcal{H}_{F},$ while the irreducible representations
of $u\left(  r\right)  $ are obtained by restricting to subspaces of states of
fixed integer particle number.

\section{One Particle Groups \label{one particle groups}}

In this section we describe the possible direct sum decompositions of the lie
algebra, $u(2s),$ of the group $U(2s)$ of unitary transformations of one
particle state space, $\mathcal{H}^{1}$ of dimension $2s,$ which lead in a
canonical fashion to families of CS's in $N$-particle state space
$\mathcal{H}^{N}$. These decompositions are based on the subgroup of special
orthogonal transformations, $SO\left(  2s,\mathbb{R}\right)  ,$ of a real
$2s$-dimensional Hilbert space, the subgroup, $U\left(  N\right)  \times
U\left(  2s-N\right)  $, formed by products of unitary transformations of
complex $N$ and $2s-N$ dimensional Hilbert spaces and the subgroup,
$USp\left(  2s\right)  $, of transformations of a complex $2s$-dimensional
Hilbert space, that are simultaneously unitary and symplectic. This latter
subgroup is isomorphic to the group of unitary transformations, $U\left(
s,\mathbb{Q}\right)  $ of a quaternionic $s$-dimensional Hilbert space. We
signify the CS's that they induce as $G/K$-CS's, where $G=U\left(  2s\right)
$ and the subgroup $K\in\left\{  SO\left(  2s,\mathbb{R}\right)  ,U\left(
N\right)  \times U\left(  2s-N\right)  ,USp\left(  2s\right)  \right\}  $.

The Lie algebra $u(2s)$ has three direct sum Cartan decompositions
\cite{Encyclo} which are in terms of the Lie algebras, $so\left(
2s,\mathbb{R}\right)  ,$ $\left[  u\left(  N\right)  \oplus u\left(
2s-N\right)  \right]  ,$ and $usp\left(  2s\right)  ,$ of its maximal compact
\cite{Encyclo} subgroups $SO\left(  2s,\mathbb{R}\right)  ,$ $\left[  U\left(
N\right)  \times U\left(  2s-N\right)  \right]  ,$ and $USp\left(  2s\right)
$, (see \cite{Gilmore_book} for the definitions of these groups). These
decompositions are%
\begin{subequations}
\begin{align}
u\left(  2s\right)   &  =so\left(  2s,\mathbb{R}\right)  \oplus so\left(
2s,\mathbb{R}\right)  ^{\bot}\label{d1}\\
u\left(  2s\right)   &  =\left[  u\left(  N\right)  \oplus u\left(
2s-N\right)  \right]  \oplus\left[  u\left(  N\right)  \oplus u\left(
2s-N\right)  \right]  ^{\bot};\quad0\leq N\leq2s\label{d2}\\
u\left(  2s\right)   &  =usp\left(  2s\right)  \oplus usp\left(  2s\right)
^{\bot} \label{d3}%
\end{align}
where the orthogonality is with respect to the killing form inner product
\cite{Encyclo} of $u(2s)$. The decomposition eqs. $\left(  \ref{d1}%
,\ref{d2},\ref{d3}\right)  $ define the following distinct coset
decompositions of $U(2s)$%
\end{subequations}
\begin{subequations}
\begin{align}
&  U\left(  2s\right)  \left/  U\left(  N\right)  \otimes U\left(
2s-N\right)  \overset{}{}\right.  ;\quad0\leq N\leq2s\label{c1}\\
&  U\left(  2s\right)  \left/  USp\left(  2s\right)  \overset{}{}\right.
\label{c2}\\
&  U\left(  2s\right)  \left/  SO\left(  2s,\mathbb{R}\right)  \overset{}%
{}\right.  \label{c3}%
\end{align}
and thus three classes of families of coherent states (see Section
\ref{cohstate}).

The motivation of using \emph{maximal} compact subgroups in eqs. $\left(
\ref{c1}\right)  ,\left(  \ref{c2}\right)  $ and $\left(  \ref{c3}\right)  $
is to produce coset spaces that are irreducible Riemannian Manifolds
\cite{Helgason78} (a Riemannian Manifold has tangent spaces that are inner
product spaces) i.e. they cannot be factored into a Riemannian product of
smaller Riemannian Manifolds. This technical feature, although important in
the classification of Riemannian Manifolds, is too restrictive when
considering manifolds of CS's, where we also need to examine decompositions
generated by subgroups of these maximal compact groups. We thus enlarge the
designation $G/K$-CS to include subgroups of $K$.

In order for any of the subgroups
\end{subequations}
\begin{equation}
\left\{  SO\left(  2s,\mathbb{R}\right)  ,\left\{  \left.  U\left(  N\right)
\otimes U\left(  2s-N\right)  \right\vert 0\leq N\leq2s\right\}  ,USp\left(
2s\right)  \right\}  ,
\end{equation}
or their subgroups, to be isotropy groups of a given $N$-electron state, the
state must carry a one dimensional unitary representation of that subgroup and
not of any larger subgroup containing that subgroup.

There are \emph{no} states in $\mathcal{H}^{N}$ that produce a one-dimensional
representation of $SO\left(  2s,\mathbb{R}\right)  $ and we conjecture (see
Appendix \ref{conjecture}) that there are \emph{no }subgroups of $SO\left(
2s,\mathbb{R}\right)  ,$ which \emph{are not} subgroups of $U\left(  N\right)
\otimes U\left(  2s-N\right)  $ or $USp\left(  2s\right)  ,$ that produce
one-dimensional representations on $\mathcal{H}^{N}$. Thus there are \emph{no
}CS manifolds in $\mathcal{H}^{N}$ based on subgroups of $SO\left(
2s,\mathbb{R}\right)  $ that cannot be described in terms of $U\left(
N\right)  \otimes U\left(  2s-N\right)  $ or $USp\left(  2s\right)  $ CS's.

The only subgroup from the set $\left\{  \left.  U\left(  N\right)  \otimes
U\left(  2s-N\right)  \right\vert 0\leq N\leq2s\right\}  $ that produces a
one-dimensional representations in $\mathcal{H}^{N}$ is $U\left(  N\right)
\otimes U\left(  2s-N\right)  .$ The $U\left(  2s\right)  \left/  U\left(
N\right)  \otimes U\left(  2s-N\right)  \overset{}{}\right.  $-coherent states
give the well known Thouless parameterization \cite{Thouless60} of the
manifold of Independent Particle States (IPS) and lead to the standard Random
Phase Approximation (RPA) for the linear response function. The subgroup
$U\left(  N\right)  \otimes U\left(  2s-N\right)  $ is actually a subgroup of
$USp\left(  2s\right)  $ when $N=2M$ and we display in section $\left(
\ref{cohstate}\right)  $ that $U\left(  2s\right)  /\left[  U\left(  N\right)
\otimes U\left(  2s-N\right)  \right]  $-CS's can in fact be described in
terms of the class of $U\left(  2s\right)  /USp\left(  2s\right)  $-CS's.

We show by construction in section $\left(  \ref{cohstate}\right)  $ that the
group $USp\left(  2s\right)  $ and its subgroups do have one dimensional
representations in $\mathcal{H}^{N}=\mathcal{H}^{2M}$ and one can form
\begin{equation}
U\left(  2s\right)  /K-\text{coherent states, where }USp\left(  2s\right)
\supseteq K, \label{USPCS}%
\end{equation}
which are manifolds of general Antisymmetrized Geminal Power (AGP) states (see
Appendix \ref{AGP}). All manifolds of CS's in $\mathcal{H}^{2M}$ based on
$U\left(  2s\right)  $ can thus be expressed in terms of the cosets eq.
$\left(  \ref{USPCS}\right)  $ which have the form
\begin{equation}
U\left(  2s\right)  \left/  \left[  USp\left(  2\omega_{1}\right)
\cdots\times USp\left(  2\omega_{p}\right)  \times SU\left(  2\right)
^{\sigma}\times SU\left(  2\xi\right)  \right]  \overset{}{}\right.  ,
\end{equation}
where
\begin{equation}
SU\left(  2\right)  ^{\sigma}=\,\overset{\sigma-\text{factors}}{\overbrace
{SU\left(  2\right)  \times\cdots\times SU\left(  2\right)  }},
\end{equation}
the index set $\left\{  \sigma,\omega_{1},\cdots,\omega_{\rho},\xi\right\}  $
characterizes the type of AGP\ reference state and its elements satisfy
\begin{align}
2\sigma+2\omega+2\xi &  =r=2s\nonumber\\
\omega &  =\sum_{1\leq j\leq\rho}\omega_{j}.
\end{align}
The index set $\left\{  \sigma,\omega_{1},\cdots,\omega_{\rho},\xi\right\}  $
is determined by the array $\left\{  \left.  c_{j}\right\vert 1\leq j\leq
s\right\}  $ of canonical canonical coefficients \cite{Weiner&ortiz04} of the
geminal that determines the AGP\ state according to:

\begin{description}
\item[$\cdot$] $\sigma=$ number of non zero, non equal $\left\{
c_{j}\right\}  $'s

\item[$\cdot$] $\omega_{j}=$ number of equal coefficients in the $j^{th}$ set

\item[$\cdot$] $\rho=$ number of sets of equal $\left\{  c_{j}\right\}  $'s

\item[$\cdot$] $\xi=$ number of zero coefficients.
\end{description}

If $\xi=0$ and $\omega=0\Leftrightarrow\rho=0,$ i.e. all the coefficients are
different and non zero, then the isotropy subgroup is
\begin{equation}
SU\left(  2\right)  ^{s}=\overset{s-\text{factors}}{\overbrace{SU\left(
2\right)  \times\cdots\times SU\left(  2\right)  }}%
\end{equation}
which is the smallest one possible, leading to the biggest coset.\ In this
case the action of coset on the reference AGP generates the manifold of all
AGP\ states. If $\xi\neq0$ the reference AGP\ is produced by a degenerate
geminal and the action of the associated coset only produces a submanifold of
AGP\ states, all of which correspond to geminals that have the same null space
$\left(  \text{see Appendix \ref{AGP}}\right)  $ as the reference AGP. If
$\rho\neq0$ then the geminal has some extreme components and again the coset
only generates a sub manifold of AGP's from the reference. If $\sigma=0$,
$\xi=0$ and $\rho=s$ then the reference state is an extreme AGP\ (see Appendix
\ref{agp}) and the action of the coset on the reference produces the manifold
of \emph{all }extreme AGP states.

All CS manifolds in $\mathcal{H}^{2M}$ based on the one particle group
$U\left(  2s\right)  $ are families of AGP\ states, some of these manifolds
are irreducible Riemannian manifolds and correspond to cosets formed by the
maximal compact subgroups $U\left(  2M\right)  \otimes U\left(  2s-2M\right)
$ and $USp\left(  2s\right)  ,$ while others are reducible and correspond to
non maximal compact subgroups $USp\left(  2\omega_{1}\right)  \cdots\times
USp\left(  2\omega_{p}\right)  \times SU\left(  2\right)  ^{\sigma}\times
SU\left(  2\xi\right)  $. Each of these manifolds of CS's have certain basic
physical properties e.g. $U\left(  2M\right)  \otimes U\left(  2s-2M\right)  $
invariant manifold describes uncorrelated Independent Particle States (IPS's),
the $USp\left(  2s\right)  $ invariant manifold describes highly correlated
extreme AGP\ states that are superconducting, while the $USp\left(
2\omega_{1}\right)  \cdots\times USp\left(  2\omega_{\rho}\right)  \times
SU\left(  2\right)  ^{\sigma}\times SU\left(  2\xi\right)  $ invariant
manifold for general $\left\{  \omega_{1},\cdots,\omega_{\rho},\sigma
,\xi\right\}  $ describe intermediate types of correlation and linear response
properties, see for a particular example \cite{Goscinski82a}, most of which
have not been explored in any depth.

\section{Coherent States\label{cohstate}}

\subsection{Definition and Construction\label{def}}

The coherent state construction is based on the Hilbert space of states
$\mathcal{H}$ that carries unitary irreducible representations $\mathcal{R}%
\left(  G\right)  $ of a compact group $G$ on $\mathcal{H}$ and a one
dimensional representation of a subgroup $K\subset G$. The groups $G$ we
consider are $U\left(  \mathcal{H}^{N}\right)  $ and $U\left(  \mathcal{H}%
^{1}\right)  $ the groups associated with $N$-electron state space and
$1$-electron state space respectively. In the finite dimensional case these
groups are the classical matrix groups and $U\left(  \mathcal{H}^{N}\right)
\equiv U\left(  \left(
%TCIMACRO{\QATOP{r}{N}}%
%BeginExpansion
\genfrac{}{}{0pt}{}{r}{N}%
%EndExpansion
\right)  \right)  $ and $U\left(  \mathcal{H}^{1}\right)  \equiv U\left(
r\right)  ,$ in general we will use the matric notation even though in an
exact treatment $r=\infty$ and one should be aware of possible infinite
dimensional pathologies. The irreducible representations we consider are
$\mathcal{R}\left(  U\left(  \left(
%TCIMACRO{\QATOP{r}{N}}%
%BeginExpansion
\genfrac{}{}{0pt}{}{r}{N}%
%EndExpansion
\right)  \right)  \right)  $ and $\mathcal{R}\left(  U\left(  r\right)
\right)  $ that are carried by $\mathcal{H}^{N},$ unless otherwise stated we
will not distinguish between the groups $G$ and $K$ and their representations
on $\mathcal{H}^{N}$.

A state $\left\vert \Phi\right\rangle \in\mathcal{H}^{N}$ is a reference state
for a group $G$ if there is an isotropy subgroup $K$ that is contained in $G$
that has $1$-dimensional representation on $\mathcal{H}^{N}$ defined by
\begin{equation}
K=\left\{  \left.  g\right\vert \quad\mathcal{R}\left(  g\right)  \left\vert
\Phi\right\rangle =e^{i\theta\left(  g\right)  }\left\vert \Phi\right\rangle
;g\in G;\theta\left(  g\right)  \in\mathbb{R}\right\}  . \label{subgpcon}%
\end{equation}
If such isotropy groups exists $G$ can be factored as%
\begin{equation}
G=\left(  G/K\right)  K.
\end{equation}


\subsection{AGP\ Coherent States}

The group $U\left(  2s\right)  $ has subgroups that leave reference
$2M$-electron AGP\ states,$\left\vert g^{\frac{N}{2}}\right\rangle $ invariant
up to an overall phase factor and thus satisfy eq. $\left(  \ref{subgpcon}%
\right)  $. AGP states are given \cite{Weiner&ortiz04} by%
\begin{equation}
\left\vert g^{\frac{N}{2}}\right\rangle =\sum_{1\leq j_{1}<j_{2}%
<\cdots<j_{\frac{N}{2}}\leq s}c_{j_{1}}\cdots c_{j_{\frac{N}{2}}}\left\vert
\eta_{j_{1}}\eta_{j_{1}+s}\cdots\eta_{j_{N}}\eta_{j_{N}+s}\right\rangle ,
\label{canexp}%
\end{equation}
where $\left\{  \left.  \left\vert \eta_{j}\right\rangle \right\vert 1\leq
j\leq r\right\}  $ is a complete orthonormal basis for $\mathcal{H}^{1},$
$\left\{  \left.  c_{j_{1}}\right\vert 1\leq j\leq s\right\}  $ are positive
real numbers called the canonical coefficients of the two electron geminal
$\left\vert g\right\rangle $ which can be expanded in canonical form as%
\begin{equation}
\left\vert g\right\rangle =\sum_{1\leq j\leq s}c_{j_{1}}\left\vert \eta
_{j}\eta_{j+s}\right\rangle . \label{gemcan}%
\end{equation}
The state $\left\vert g^{\frac{N}{2}}\right\rangle $ can also be expressed in
terms of the coset $SO\left(  2r,\mathbb{R}\right)  /SU\left(  r\right)  $ of
the group $SO\left(  2r,\mathbb{R}\right)  $ acting on the zero particle
vacuum state $\left\vert \phi\right\rangle $ of $\mathcal{H}_{F}$
\cite{Deumens87} as
\begin{equation}
\left\vert g^{\frac{N}{2}}\right\rangle =P_{\frac{N}{2}}\exp\left\{
\sum_{1\leq j\leq s}c_{j_{1}}b_{j}^{\dagger}b_{j+s}^{\dagger}\right\}
\left\vert \phi\right\rangle
\end{equation}
where $P_{\frac{N}{2}}$ is the orthogonal projector onto $\mathcal{H}^{N}$ and
$\left\{  \left.  b_{j}^{\dagger}\right\vert 1\leq j\leq s\right\}  $ are the
creators for the one particle states $\left\{  \left.  \left\vert \eta
_{j}\right\rangle \right\vert 1\leq j\leq r\right\}  $ i.e.
\begin{equation}
b_{j}^{\dagger}\left\vert \phi\right\rangle =\left\vert \eta_{j}\right\rangle
.
\end{equation}


\subsection{Coset Factorization for $U\left(  2s\right)  $\label{factor}}

In general the action of $U\left(  2s\right)  $ on a reference AGP,
$\left\vert g^{\frac{N}{2}}\right\rangle $, state produces a manifold,
$\mathcal{M}\left(  \sigma,\omega_{1},\cdots,\omega_{\rho},\xi\right)  $ of
AGP states given by
\begin{equation}
\left\{  \left[  U\left(  2s\right)  \left/  K\overset{}{}\right.  \right]
\times K\right\}  \left\vert g^{\frac{N}{2}}\right\rangle . \label{su2act}%
\end{equation}
where
\begin{equation}
K=USp\left(  2\omega_{1}\right)  \cdots\times USp\left(  2\omega_{p}\right)
\times SU\left(  2\right)  ^{\sigma}\times SU\left(  2\xi\right)  .
\end{equation}
Each element of the manifold can be expressed as
\begin{equation}
\left\vert g\left(  c\right)  ^{\frac{N}{2}}\right\rangle =ck_{\omega_{1}%
}\cdots k_{\omega_{\rho}}k_{\sigma}k_{\xi}\left\vert g^{\frac{N}{2}%
}\right\rangle =c\left\vert g^{\frac{N}{2}}\right\rangle ,
\end{equation}
where the isotropy subgroup is given by the products $k_{\omega_{1}}\cdots
k_{\omega_{\rho}}k_{\sigma}k_{\xi}$ where%
\begin{equation}
k_{\omega_{j}}\in USp\left(  2\omega_{j}\right)  ;1\leq j\leq\rho,k_{\sigma
}\in SU\left(  2\right)  ^{\sigma},k_{\xi}\in SU\left(  2\xi\right)  ,
\end{equation}
$c\in U\left(  2s\right)  /K$ and the reference AGP is generated by a geminal
$\left\vert g\right\rangle $ in the class $\left\{  \rho,\sigma,\xi\right\}
.$ The most general AGP manifold is produced when $\sigma=s$ ($\Leftrightarrow
\left(  \rho=0\text{ \& }\xi=0\right)  $ ) so it is advantages to first
consider that case.

\subsection{$SU\left(  2\right)  ^{s}$ factors}

The isotropy sub group, $K$, for an AGP reference state that has non equal and
nonzero canonical coefficients $\left\{  \left.  c_{j}\right\vert 1\leq j\leq
s\right\}  ,$ i.e. $\sigma=s,$ is the product group $SU\left(  2\right)
^{s},$ which produces the coset decomposition
\begin{equation}
U\left(  2s\right)  =\left[  U\left(  2s\right)  /SU\left(  2\right)
^{s}\right]  \times SU\left(  2\right)  ^{s}.
\end{equation}
One can express coset representatives in a factored form as the product
$G_{+1}G_{0}G_{-1}$, \cite{Zhang90} where $\left\{  G_{+1},G_{0}%
,G_{-1}\right\}  $ are subgroups of the linear group $GL\left(  2s,\mathbb{C}%
\right)  ,\ $which is the complexification of $U\left(  2s\right)  ,$ and the
group $G_{-1}$ leaves the reference AGP state invariant (note that it is
different to the real compact isotropy groups). Restricting attention to the
connected components of these subgroups, that contain the identity element,
one can express this factorization in terms of Lie algebras as
\begin{equation}
U\left(  2s\right)  =e^{N_{+}}e^{N_{0}}e^{N_{-}}e^{\mathcal{K}}%
\end{equation}
where
\begin{equation}
\mathcal{K}=\,\overset{s\text{-subspaces}}{\overbrace{su\left(  2\right)
\oplus\cdots\oplus su\left(  2\right)  }},
\end{equation}
$\left\{  N_{+1},N_{0},N_{-1}\right\}  $ are Lie subalgebras of $gl\left(
2s,\mathbb{C}\right)  ,$ the Lie algebra of $GL\left(  2s,\mathbb{C}\right)
$. The subalgebras $\left\{  N_{-1},\mathcal{K}\right\}  $ annihilate
$\left\vert g^{\frac{N}{2}}\right\rangle $ i.e.
\begin{equation}
b\left\vert g^{\frac{N}{2}}\right\rangle =0\quad\forall b\in N_{-1}%
\oplus\mathcal{K},
\end{equation}
while the subalgebras $\left\{  N_{+1},N_{0}\right\}  $ have the property
\begin{equation}
b\left\vert g^{\frac{N}{2}}\right\rangle \neq0\quad\forall b\in N_{+1}\oplus
N_{0}.
\end{equation}
Introducing coordinate systems $\left\{  \mathbf{z}_{+},\mathbf{\lambda,z}%
_{-},\mathbf{\eta}\right\}  $ with respect to fixed bases of $\left\{
N_{+1},N_{0},N_{-1},\mathcal{K}\right\}  ,$ that depend on the reference AGP
state, one can parameterize group elements $u\in U\left(  2s\right)  $ as
\begin{equation}
u\left(  \mathbf{\tilde{z},\eta}\right)  =\tilde{c}\left(  \mathbf{\tilde{z}%
}\right)  k\left(  \mathbf{\eta}\right)  =c_{+1}\left(  \mathbf{z}_{+}\right)
c_{0}\left(  \mathbf{\lambda}\right)  c_{-1}\left(  \mathbf{z}_{-}\right)
k\left(  \mathbf{\eta}\right)  .
\end{equation}
The action of $U\left(  2s\right)  $ on the AGP state is
\begin{equation}
u\left(  \mathbf{\tilde{z},\eta}\right)  \left\vert g^{\frac{N}{2}%
}\right\rangle =\tilde{c}\left(  \mathbf{\tilde{z}}\right)  k\left(
\mathbf{\eta}\right)  \left\vert g^{\frac{N}{2}}\right\rangle =c_{+1}\left(
\mathbf{z}_{+}\right)  c_{0}\left(  \mathbf{\lambda}\right)  c_{-1}\left(
\mathbf{z}_{-}\right)  k\left(  \mathbf{\eta}\right)  \left\vert g^{\frac
{N}{2}}\right\rangle =c_{+1}\left(  \mathbf{z}_{+}\right)  c_{0}\left(
\mathbf{\lambda}\right)  \left\vert g^{\frac{N}{2}}\right\rangle ,
\end{equation}
showing that the non redundant coordinates $\mathbf{z}\equiv\left(
\mathbf{z}_{+},\mathbf{\lambda}\right)  $ describe a configuration space, that
is a sub manifold $\mathcal{M}$ (in general nonlinear) of $\mathcal{H}^{N}$
with a coordinate system $\left\{  \mathbf{z}\right\}  ,$ that parameterize
the set of $\left\{  \left\vert \Phi\left(  \mathbf{z}\right)  \right\rangle
\right\}  $ of $U\left(  2s\right)  /SU\left(  2\right)  ^{s}$-coherent
states
\begin{equation}
\left\vert \Phi\left(  \mathbf{z}\right)  \right\rangle =c_{+1}\left(
\mathbf{z}_{+}\right)  c_{0}\left(  \mathbf{\lambda}\right)  \left\vert
g^{\frac{N}{2}}\right\rangle \equiv c\left(  \mathbf{z}\right)  \left\vert
g^{\frac{N}{2}}\right\rangle ,
\end{equation}
and at the same time represents the coset space $U\left(  2s\right)
/SU\left(  2\right)  ^{s}$. The reference AGP is generated by a geminal
$\left\vert g\right\rangle $ in the class $\left\{  \rho,\sigma,\xi\right\}
=\left\{  0,s,0\right\}  $

The AGP\ state dependent bases one chooses are
\begin{subequations}
\begin{align}
N_{+1}  &  =\mathbb{C-}\operatorname{linspan}\left\{  c_{j}a_{k}^{\dagger
}a_{j}-\sigma\left(  j\right)  \sigma\left(  k\right)  c_{k}a_{-j}^{\dagger
}a_{-k};\right. \nonumber\\
&  \quad\quad\quad\quad\quad\quad\quad\quad\quad\quad\quad\quad\quad\quad
\quad\quad\quad\left.  -s\leq j<k\leq s,\left\vert j\right\vert \neq\left\vert
k\right\vert ;\,\sigma\left(  j\right)  =\operatorname{sign}\left(  j\right)
\right\} \label{l1}\\
N_{0}  &  =\mathbb{C-}\operatorname{linspan}\left\{  \left(  a_{j}^{\dagger
}a_{j}+a_{-j}^{\dagger}a_{-j}\right)  ,\left(  a_{j}^{\dagger}a_{-j}%
+a_{-j}^{\dagger}a_{j}\right)  ,\left(  a_{j}^{\dagger}a_{j}-a_{-j}^{\dagger
}a_{-j}\right)  ,\right. \nonumber\\
&  \quad\quad\quad\quad\quad\quad\quad\quad\quad\quad\quad\quad\quad\quad
\quad\quad\quad\quad\quad\quad\quad\quad\quad\left.  i\left(  a_{j}^{\dagger
}a_{-j}-a_{-j}^{\dagger}a_{j}\right)  ;1\leq j\leq s\right\} \label{l2}\\
N_{-1}  &  =\mathbb{C-}\operatorname{linspan}\left\{  c_{k}a_{k}^{\dagger
}a_{j}+\sigma\left(  j\right)  \sigma\left(  k\right)  c_{j}a_{-j}^{\dagger
}a_{-k}\right. \nonumber\\
&  \quad\quad\quad\quad\quad\quad\quad\quad\quad\quad\quad\quad\quad\quad
\quad\quad\quad\left.  -s\leq j<k\leq s,\left\vert j\right\vert \neq\left\vert
k\right\vert ;\,\sigma\left(  j\right)  =\operatorname{sign}\left(  j\right)
\right\} \label{l3}\\
\mathcal{K}  &  =\mathbb{R-}\operatorname{linspan}\left\{  i\left(
a_{j}^{\dagger}a_{-j}+a_{-j}^{\dagger}a_{j}\right)  ,i\left(  a_{j}^{\dagger
}a_{j}-a_{-j}^{\dagger}a_{-j}\right)  ,\left(  a_{j}^{\dagger}a_{-j}%
-a_{-j}^{\dagger}a_{j}\right)  ;\quad1\leq j\leq s\right\}  , \label{l4}%
\end{align}
where we have translated the indexing of the CGSO's by
\end{subequations}
\begin{equation}
\left\vert \eta_{j}\right\rangle \rightarrow\left\vert \eta_{j-s}\right\rangle
\quad1\leq j\leq2s.
\end{equation}
It is useful to note that
\begin{equation}
N_{0}=\left\{  \mathcal{K\oplus N}\right\}  ^{\mathbb{C}}%
\end{equation}
where
\begin{equation}
\mathcal{N=\,}%
%TCIMACRO{\tbigoplus _{1\leq j\leq s}}%
%BeginExpansion
{\textstyle\bigoplus_{1\leq j\leq s}}
%EndExpansion
\left\{  \mathbb{R-}\operatorname{linspan}\left\{  \left.  \left(
a_{j}^{\dagger}a_{j}+a_{-j}^{\dagger}a_{-j}\right)  \right\vert \quad1\leq
j\leq s\right\}  \right\}  .
\end{equation}
and $\left\{  {}\right\}  ^{\mathbb{C}}$ denotes complexification. The factors
in the group decomposition are%
\begin{align}
c_{+1}\left(  \mathbf{z}_{+}\right)   &  =\exp\left\{  \sum
\limits_{\substack{-s\leq j<k\leq s\\\left\vert j\right\vert \neq\left\vert
k\right\vert }}z_{+jk}\left\{  c_{j}a_{k}^{\dagger}a_{j}-\sigma\left(
j\right)  \sigma\left(  k\right)  c_{k}a_{-j}^{\dagger}a_{-k}\right\}
\right\}  =\exp\sum\limits_{_{\substack{-s\leq j<k\leq s\\\left\vert
j\right\vert \neq\left\vert k\right\vert }}}z_{+jk}q_{jk}^{\dagger
};\label{co1}\\
c_{-1}\left(  \mathbf{z}_{-}\right)   &  =\exp\left\{  \sum
\limits_{\substack{-s\leq j<k\leq s\\\left\vert j\right\vert \neq\left\vert
k\right\vert }}z_{-jk}\left\{  c_{k}a_{k}^{\dagger}a_{j}+\sigma\left(
j\right)  \sigma\left(  k\right)  c_{j}a_{-j}^{\dagger}a_{-k}\right\}
\right\}  =\exp\sum\limits_{_{\substack{-s\leq j<k\leq s\\\left\vert
j\right\vert \neq\left\vert k\right\vert }}}z_{-jk}q_{jk}\label{co2}\\
c_{0}\left(  \mathbf{\lambda}\right)   &  =\exp\left\{  \sum
\limits_{_{\substack{1\leq j<\leq s\\1\leq k\leq4}}}\lambda_{jk}n_{jk}\right\}
\label{co3}\\
h\left(  \mathbf{\eta}\right)   &  =\exp\left\{  \sum
\limits_{_{\substack{1\leq j<\leq s\\1\leq k\leq3}}}\eta_{jk}u_{jk}\right\}
\label{co4}%
\end{align}
where $\left\{  z_{+jk},z_{-jk}\right\}  \in\mathbb{C},\left\{  \lambda
_{jk},\eta_{jk}\right\}  \in\mathbb{R}$ , and $\left\{  n_{jk},u_{jk}\right\}
$ are given by equ's. $\left(  \ref{l2}\right)  $ and $\left(  \ref{l4}%
\right)  $ respectively.

In order that the transformation $c_{+1}\left(  \mathbf{z}_{+}\right)
c_{0}\left(  \mathbf{\lambda}\right)  c_{-1}\left(  \mathbf{z}_{-}\right)
k\left(  \mathbf{\eta}\right)  $ be unitary the values of $\left\{
\mathbf{z}_{+},\mathbf{\lambda,z}_{-}\right\}  $ must be constrained resulting
in dependencies between the variables. The action of such unitary
transformations on a normalized reference state produce normalized coherent
states that can be expressed in terms of the coset representatives
$c_{+1}\left(  \mathbf{z}_{+}\right)  c_{0}\left(  \mathbf{\lambda}\right)  $.
However there are dependencies between the variables $\left\{  \mathbf{z}%
_{+}\right\}  $ and $\left\{  \mathbf{\lambda}\right\}  $ which are often
convenient to avoid. One method of achieving this is to consider unnormalized
coherent states by modifying the coset representative $c_{0}\left(
\mathbf{\lambda}\right)  $ and explicitly considering the norm of the state in
various expectation value expressions.

\subsection{$USp\left(  \omega_{j}\right)  $ factors}

The changes from the $SU\left(  2\right)  ^{s}$ case can be illustrated by an
$\sigma=s-2,\rho=1,\omega_{1}=2$ example, which clearly leads to the general
picture for $\rho>1$. Here we have $s-2$ eigenvalues, $\left\{  \left.
c_{j}^{2}\right\vert 3\leq j\leq s\right\}  $ of $gg^{\dag}$ that are doubly
degenerate and one fourfold degenerate eigenvalue, $c_{1}^{2}$, of $gg^{\dag}%
$, which is associated with two doubly degenerate eigenvalues, $\left\{
c_{1},-c_{1}\right\}  $ of $g$ affiliated with CGSO's that are linear
combinations of $\left\{  \left\vert \eta_{1}\right\rangle ,\left\vert
\eta_{2}\right\rangle \right\}  $ and $\left\{  \left\vert \eta_{-1}%
\right\rangle ,\left\vert \eta_{-2}\right\rangle \right\}  $ respectively
\cite{Weiner&ortiz04}. In this case unitary transformations, $u\left(
\mathbf{\tilde{z},\eta}\right)  $, belonging to $U\left(  2s\right)  $ can be
expressed in factored form as
\begin{align}
u\left(  \mathbf{\tilde{z},\eta}\right)   &  =c\left(  \mathbf{\tilde{z}%
}\right)  k_{\sigma}\left(  \mathbf{\eta}_{\sigma}\right)  k_{\omega}\left(
\mathbf{\eta}_{\omega}\right) \nonumber\\
&  =c_{+1}\left(  \mathbf{z}_{\sigma+}\right)  u_{+1}\left(  \mathbf{z}%
_{\omega+}\right)  u_{0}\left(  \mathbf{\lambda}_{\omega}\right)  c_{0}\left(
\mathbf{\lambda}_{\sigma}\right)  u_{-1}\left(  \mathbf{z}_{\omega-}\right)
c_{-1}\left(  \mathbf{z}_{-}\right)  k_{\sigma}\left(  \mathbf{\eta}_{\sigma
}\right)  k_{\omega}\left(  \mathbf{\eta}_{\omega}\right)  . \label{gpprod}%
\end{align}
where $k_{\sigma}\left(  \mathbf{\eta}_{\sigma}\right)  \in SU\left(
2\right)  ^{s-2}$ and $c_{\sigma0}\left(  \mathbf{\lambda}_{\sigma}\right)
\in\exp N_{0}$ and are expanded in the basis eqs.$\left(  \ref{l2}\right)  $
and $\left(  \ref{l4}\right)  $ respectively with the operators referring to
CGSO's $\left\{  \left.  \left\vert \eta_{j}\right\rangle \right\vert
j=\pm1,\pm2\right\}  $ deleted. The operators $c_{+1}\left(  \mathbf{z}%
_{\sigma+}\right)  \in\exp N_{+1}$ and $c_{-1}\left(  \mathbf{z}_{\sigma
-}\right)  \in\exp N_{-1}$ are expanded in the basis eqs.$\left(
\ref{l1}\right)  $ and $\left(  \ref{l3}\right)  $ with the operators
referring to the pairs $\left\{  \left.  \left(  \left\vert \eta
_{j}\right\rangle ,\left\vert \eta_{k}\right\rangle \right)  \right\vert
j=\pm1,k=\pm1\right\}  $ of CGSO's deleted. The factors $u_{+1}\left(
\mathbf{z}_{\omega+}\right)  \in\exp U_{+11},u_{-1}\left(  \mathbf{z}%
_{\omega-}\right)  \in\exp U_{-11},u_{\omega0}\left(  \mathbf{\lambda}%
_{\omega}\right)  \in\exp U_{01},k_{\omega}\left(  \mathbf{\eta}_{\omega
}\right)  \in USp\left(  4\right)  $, where $U_{+1},U_{0},U_{-1}$ and
$usp\left(  4\right)  $ are expanded in the following bases%

\begin{subequations}
\begin{align}
U_{+11}  &  =\mathbb{C-}\operatorname{linspan}\left\{  a_{k}^{\dagger}%
a_{j}+\sigma\left(  j\right)  \sigma\left(  k\right)  a_{-j}^{\dagger}%
a_{-k};\quad j=1,-1;k=2,-2\right\} \label{u1}\\
U_{01}  &  =\mathbb{C-}\operatorname{linspan}\left\{  \left(  a_{j}^{\dagger
}a_{j}+a_{-j}^{\dagger}a_{-j}\right)  ,\left(  a_{j}^{\dagger}a_{j}%
-a_{-j}^{\dagger}a_{-j}\right)  ,\left(  a_{j}^{\dagger}a_{-j}+a_{-j}%
^{\dagger}a_{j}\right)  ,\right. \nonumber\\
&  i\left(  a_{j}^{\dagger}a_{-j}-a_{-j}^{\dagger}a_{j}\right)  ,i\left(
a_{k}^{\dagger}a_{j}-\sigma\left(  j\right)  \sigma\left(  k\right)
a_{-j}^{\dagger}a_{-k}-a_{j}^{\dagger}a_{k}+\sigma\left(  j\right)
\sigma\left(  k\right)  a_{-k}^{\dagger}a_{-j}\right)  ,\nonumber\\
&  \left.  \left(  a_{k}^{\dagger}a_{j}-\sigma\left(  j\right)  \sigma\left(
k\right)  a_{-j}^{\dagger}a_{-k}+a_{j}^{\dagger}a_{k}-\sigma\left(  j\right)
\sigma\left(  k\right)  a_{-k}^{\dagger}a_{-j}\right)  ;\quad
j=1,-1;k=2,-2\right\} \label{u2}\\
U_{-11}  &  =\mathbb{C-}\operatorname{linspan}\left\{  a_{k}^{\dagger}%
a_{j}-\sigma\left(  j\right)  \sigma\left(  k\right)  a_{-j}^{\dagger}%
a_{-k};\quad j=1,-1;k=2,-2\right\} \label{u3}\\
usp\left(  2\omega_{1}\right)   &  =usp\left(  4\right)  =\mathbb{R-}%
\operatorname{linspan}\left\{  i\left(  a_{j}^{\dagger}a_{j}-a_{-j}^{\dagger
}a_{-j}\right)  ,i\left(  a_{j}^{\dagger}a_{-j}+a_{-j}^{\dagger}a_{j}\right)
,\left(  a_{j}^{\dagger}a_{-j}-a_{-j}^{\dagger}a_{j}\right)  ,\right.
\nonumber\\
&  \left(  a_{k}^{\dagger}a_{j}-\sigma\left(  j\right)  \sigma\left(
k\right)  a_{-j}^{\dagger}a_{-k}-a_{j}^{\dagger}a_{k}+\sigma\left(  j\right)
\sigma\left(  k\right)  a_{-k}^{\dagger}a_{-j}\right)  ,\nonumber\\
&  \left.  i\left(  a_{k}^{\dagger}a_{j}-\sigma\left(  j\right)  \sigma\left(
k\right)  a_{-j}^{\dagger}a_{-k}+a_{j}^{\dagger}a_{k}-\sigma\left(  j\right)
\sigma\left(  k\right)  a_{-k}^{\dagger}a_{-j}\right)  ;\quad
j=1,-1;k=2,-2\right\}  , \label{u4}%
\end{align}
where $\sigma\left(  j\right)  =\operatorname{sign}\left(  j\right)  .$ The
manifold of $U\left(  2s\right)  /\left[  USp\left(  4\right)  \times
SU\left(  2\right)  ^{s-1}\right]  $-coherent states is then given by
\end{subequations}
\begin{equation}
\left\vert \Phi\left(  \mathbf{z}\right)  \right\rangle =c_{+1}\left(
\mathbf{z}_{+\sigma}\right)  u_{+1}\left(  \mathbf{z}_{+\omega}\right)
u_{01}\left(  \mathbf{\lambda}_{\omega_{1}}\right)  c_{0}\left(
\mathbf{\lambda}_{\sigma}\right)  \left\vert g^{\frac{N}{2}}\right\rangle ,
\end{equation}
where $\mathbf{z\equiv}\left(  \mathbf{z}_{+\sigma},\mathbf{z}_{+\omega
},\mathbf{\lambda}_{\omega},\mathbf{\lambda}_{\sigma}\right)  $ and the
reference AGP is generated by a geminal $\left\vert g\right\rangle $ in the
class $\left\{  \rho,\sigma,\xi\right\}  =\left\{  1,s-1,0\right\}  .$

If $\rho>2$ then one deletes the appropriate operators in the bases $\left\{
N_{+1},N_{0},N_{-1},\mathcal{K}\right\}  $ defined in eqs. $\left(
\ref{l1}\right)  -\left(  \ref{l4}\right)  $ and defines new bases $\left\{
\left.  \left\{  U_{+1j},U_{0j},U_{-1j},usp\left(  2\omega_{j}\right)
\right\}  \right\vert 1\leq j\leq\rho\right\}  $ according to eqs. $\left(
\ref{u1}\right)  -\left(  \ref{u4}\right)  .$ The subalgebras $\left\{
\left.  U_{-1j},usp\left(  2\omega_{j}\right)  \right\vert 1\leq j\leq
\rho\right\}  $ annihilate $\left\vert g^{\frac{N}{2}}\right\rangle $ i.e.
\begin{equation}
b\left\vert g^{\frac{N}{2}}\right\rangle =0\quad\forall b\in\bigoplus
\limits_{1\leq j\leq\rho}U_{-1j}\oplus\bigoplus\limits_{1\leq j\leq\rho
}usp\left(  2\omega_{j}\right)  ,
\end{equation}
while the subalgebras $\left\{  \left.  U_{+1j},U_{0j}\right\vert 1\leq
j\leq\rho\right\}  $ have the property
\begin{equation}
b\left\vert g^{\frac{N}{2}}\right\rangle \neq0\quad\forall b\in\bigoplus
\limits_{1\leq j\leq\rho}U_{+1j}\oplus\bigoplus\limits_{1\leq j\leq\rho}%
U_{0j}.
\end{equation}
The $U\left(  2s\right)  /\left[
%TCIMACRO{\tprod _{1\leq j\leq\rho}}%
%BeginExpansion
{\textstyle\prod_{1\leq j\leq\rho}}
%EndExpansion
USp\left(  2\omega_{j}\right)  \times SU\left(  2\right)  ^{s-\omega}\right]
$-coherent state manifold is then given by
\begin{equation}
\left\vert \Phi\left(  \mathbf{z}\right)  \right\rangle =c_{+1}\left(
\mathbf{z}_{\sigma+}\right)
%TCIMACRO{\tprod _{1\leq j\leq\rho}}%
%BeginExpansion
{\textstyle\prod_{1\leq j\leq\rho}}
%EndExpansion
\left\{  u_{+1j}\left(  \mathbf{z}_{+j}\right)  u_{0j}\left(  \mathbf{\lambda
}_{j}\right)  \right\}  c_{0}\left(  \mathbf{\lambda}_{\sigma}\right)
\left\vert g^{\frac{N}{2}}\right\rangle .
\end{equation}


The manifold of extreme AGP's $\mathcal{M}\left(  0,s,0\right)  $ are
$U\left(  2s\right)  /USp\left(  2s\right)  $-coherent state, which have been
discussed in super conductivity \cite{Yukalov00} is given by
\begin{equation}
\left\vert \Phi\left(  \mathbf{z}\right)  \right\rangle =u_{+11}\left(
\mathbf{z}_{+1}\right)  u_{01}\left(  \mathbf{\lambda}_{1}\right)  \left\vert
g^{\frac{N}{2}}\right\rangle .
\end{equation}


\subsection{$SU\left(  2\xi\right)  $ factors}

The changes from $SU\left(  2\right)  ^{s}$ can this time be illustrated by an
$\sigma=s-1,\rho=0,\xi=1$ example, which also clearly leads to the general
picture for $\xi>1$. Here we have $s-1$ eigenvalues, $\left\{  \left.
c_{j}^{2}\right\vert 2\leq j\leq s\right\}  $ of $gg^{\dag}$ that are doubly
degenerate and a doubly degenerate zero eigenvalue of $gg^{\dag}$, which is
associated with the CGSO's $\left\{  \left\vert \eta_{-1}\right\rangle
,\left\vert \eta_{1}\right\rangle \right\}  .$ In this case unitary
transformations, $u\left(  \mathbf{\tilde{z},\eta}\right)  $, belonging to
$U\left(  2s\right)  $ can be expressed in factored form as
\begin{equation}
u\left(  \mathbf{\tilde{z},\eta}\right)  =c\left(  \mathbf{\tilde{z}}\right)
k_{\sigma}\left(  \mathbf{\eta}_{\sigma}\right)  k_{\xi}\left(  \mathbf{\eta
}_{\xi}\right)  =c_{+1}\left(  \mathbf{z}_{\sigma+}\right)  v_{0}\left(
\mathbf{\lambda}_{\xi}\right)  c_{0}\left(  \mathbf{\lambda}_{\sigma}\right)
c_{-1}\left(  \mathbf{z}_{-}\right)  k_{\sigma}\left(  \mathbf{\eta}_{\sigma
}\right)  k_{\xi}\left(  \mathbf{\eta}_{\xi}\right)  , \label{untrans}%
\end{equation}
where $k_{\sigma}\left(  \mathbf{\eta}_{\sigma}\right)  \in SU\left(
2\right)  ^{s-1}$ and $c_{0}\left(  \mathbf{\lambda}_{\sigma}\right)  \in\exp
N_{0}$ which are expanded in the basis eqs.$\left(  \ref{l2}\right)  $,
$\left(  \ref{l4}\right)  $ with the operators referring to CGSO's $\left\{
\left.  \left\vert \eta_{j}\right\rangle \right\vert j=\pm1\right\}  $
deleted. The definitions of $c_{+1}\left(  \mathbf{z}_{\sigma+}\right)
\in\exp N_{+1}$ and $c_{-1}\left(  \mathbf{z}_{\sigma-}\right)  \in\exp
N_{-1}$ are unchanged and are still expanded in the basis eqs.$\left(
\ref{l1}\right)  $ and $\left(  \ref{l3}\right)  $. The factors $v_{0}\left(
\mathbf{\lambda}_{\xi}\right)  \in\exp V_{0},k_{\xi}\left(  \mathbf{\eta}%
_{\xi}\right)  \in SU\left(  2\right)  $, where $su\left(  2\right)  $ is
expanded in the following bases%
\begin{equation}
su\left(  2\right)  =\mathbb{R-}\operatorname{linspan}\left\{  i\left(
a_{1}^{\dagger}a_{-1}+a_{-1}^{\dagger}a_{1}\right)  ,i\left(  a_{1}^{\dagger
}a_{1}-a_{-1}^{\dagger}a_{-1}\right)  ,\left(  a_{1}^{\dagger}a_{-1}%
-a_{-1}^{\dagger}a_{1}\right)  \right\}
\end{equation}
and $V_{0}$ is given by
\begin{equation}
\left\{  su\left(  2\right)  \mathcal{\oplus\,N}\right\}  ^{\mathbb{C}}%
\end{equation}
where
\begin{equation}
\mathcal{N}=\mathbb{R-}\operatorname{linspan}\left\{  \left(  a_{1}^{\dagger
}a_{1}+a_{-1}^{\dagger}a_{-1}\right)  \right\}  .
\end{equation}
The subalgebra $\left\{  su\left(  2\xi\right)  \right\}  $ annihilates
$\left\vert g^{\frac{N}{2}}\right\rangle $ i.e.
\begin{equation}
b\left\vert g^{\frac{N}{2}}\right\rangle =0\quad\forall b\in su\left(
2\xi\right)  .
\end{equation}
As the class of the reference AGP\ changes from $\left(  s-1,0,1\right)  $
to$\left(  s-p,0,p\right)  $ it is clear how the definitions of the factors
eq. $\left(  \ref{untrans}\right)  $ evolve. It is easy to confirm that the
$\left(  M,0,s-M\right)  $ cannot in fact exist, while the class $\left(
0,M,s-M\right)  $ does and such reference AGP's are IPS's and the $U\left(
2s\right)  /\left[  USp\left(  2M\right)  \times SU\left(  2s-2M\right)
\right]  $-coherent state manifold is in fact the Thouless $U\left(
2s\right)  /\left[  U\left(  2M\right)  \times U\left(  2s-2M\right)  \right]
$-coherent state manifold.

As an illustration of the preceding forms of the cosets of $U\left(
\mathcal{H}^{1}\right)  $ we can consider their action on an AGP\ state
$\left\vert g^{\frac{N}{2}}\right\rangle $ determined by a geminal
\begin{equation}
\left\vert g\right\rangle =\sum_{1\leq j\leq s}c_{j_{1}}\left\vert \eta
_{j}\eta_{j+s}\right\rangle
\end{equation}
given by
\begin{equation}
c_{+1}\left(  \mathbf{z}_{+}\right)  c_{0}\left(  \mathbf{\lambda}\right)
c_{+1}\left(  \mathbf{z}_{\sigma+}\right)  v_{0}\left(  \mathbf{\lambda}_{\xi
}\right)
%TCIMACRO{\tprod _{1\leq j\leq\rho}}%
%BeginExpansion
{\textstyle\prod_{1\leq j\leq\rho}}
%EndExpansion
\left\{  u_{+1j}\left(  \mathbf{z}_{+j}\right)  u_{0j}\left(  \mathbf{\lambda
}_{j}\right)  \right\}  c_{0}\left(  \mathbf{\lambda}_{\sigma}\right)
\left\vert g^{\frac{N}{2}}\right\rangle
\end{equation}
$\sigma=$ number of non zero, non equal $\left\{  c_{j}\right\}  $'s

\begin{description}
\item[$\cdot$] $\omega_{j}=$ number of equal coefficients in the $j^{th}$ set

\item[$\cdot$] $\rho=$ number of sets of equal $\left\{  c_{j}\right\}  $'s

\item[$\cdot$] $\xi=$ number of zero coefficients.
\end{description}

\subsection{$U\left(  \mathcal{H}^{N}\right)  /U\left(  1\right)  $-Coherent
States}

In order to examine the properties of general state excitation operators in
terms of CS's it is necessary and illuminating to consider the \emph{exact
}quantum evolution equations expressed in terms of a CS manifold. The
resulting classical equations of motion defined on such a manifold are
equivalent to the full TDSE. The configuration space formed by these coherent
sates represents all of the states $\left\vert \Psi\right\rangle
\in\mathcal{H}^{N}$ via a parameterization $\left\{  \left\vert \Phi\left(
\mathbf{z}\right)  \right\rangle \right\}  $. The group $G$ in this case is
the full unitary group $U\left(  \left(
%TCIMACRO{\QATOP{r}{N}}%
%BeginExpansion
\genfrac{}{}{0pt}{}{r}{N}%
%EndExpansion
\right)  \right)  $ of $N$-electron state space and the isotropy subgroup is
\begin{equation}
K=\left\{  \left.  g\right\vert g\in G,g\left\vert \Phi_{1}\right\rangle
=e^{i\theta\left(  g\right)  }\left\vert \Phi_{1}\right\rangle \right\}
=\left\{  \exp\left(  i\lambda\left(  g\right)  \left\vert \Phi_{1}%
\right\rangle \left\langle \Phi_{1}\right\vert \right)  \exp\left(
i\sum_{i,j\neq1}\lambda_{ij}\left(  g\right)  \left\vert \Phi_{i}\right\rangle
\left\langle \Phi_{j}\right\vert \right)  \right\}
\end{equation}
The coset representatives are
\begin{equation}
G/K=\left\{  \exp\left(  \sum_{i>1}\eta_{i}\left\vert \Phi_{i}\right\rangle
\left\langle \Phi_{1}\right\vert \right)  \exp\left(  \eta_{1}\left\vert
\Phi_{1}\right\rangle \left\langle \Phi_{1}\right\vert \right)  \exp\left(
-\sum_{i>1}\bar{\eta}_{i}\left\vert \Phi_{i}\right\rangle \left\langle
\Phi_{1}\right\vert \right)  \right\}
\end{equation}
where $\left\{  \left.  \left\vert \Phi_{i}\right\rangle \right\vert
\quad1\leq i\leq\left(
%TCIMACRO{\QATOP{r}{N}}%
%BeginExpansion
\genfrac{}{}{0pt}{}{r}{N}%
%EndExpansion
\right)  \right\}  $ is a complete orthonormal basis for $\mathcal{H}^{N}$.
The family of CS's associated with this coset is defined as
\begin{equation}
\left\vert \Phi\left(  \mathbf{\eta}\right)  \right\rangle =\left\{
\exp\left(  \sum_{i>1}\eta_{i}\left\vert \Phi_{i}\right\rangle \left\langle
\Phi_{1}\right\vert \right)  \exp\left(  \eta_{1}\left\vert \Phi
_{i}\right\rangle \left\langle \Phi_{1}\right\vert \right)  \right\}
\left\vert \Phi_{1}\right\rangle .
\end{equation}
It is convenient to introduce a new equivalent coordinate system $\left\{
\mathbf{z}\right\}  $ by setting
\begin{align}
\eta_{1}  &  =\ln\left(  z_{1}\right)  \Leftrightarrow z_{1}=e^{\eta_{1}}\\
\eta_{j}  &  =\frac{z_{j}}{z_{1}};\quad2\leq j\leq\left(
%TCIMACRO{\QATOP{r}{N}}%
%BeginExpansion
\genfrac{}{}{0pt}{}{r}{N}%
%EndExpansion
\right)  \Leftrightarrow z_{j}=\eta_{j}e^{\eta_{1}}%
\end{align}
which leads to the following form for the coherent states
\begin{equation}
\left\vert \Phi\left(  \mathbf{z}\right)  \right\rangle =\sum_{1\leq
j\leq\left(
%TCIMACRO{\QATOP{r}{N}}%
%BeginExpansion
\genfrac{}{}{0pt}{}{r}{N}%
%EndExpansion
\right)  }z_{j}\left\vert \Phi_{j}\right\rangle \label{CI}%
\end{equation}
Clearly eq. $\left(  \ref{CI}\right)  $ is the familiar CI expansion of a
general $N$-fermion state and the parameters $\left\{  z_{j}\right\}  $ the
configurational interaction coefficients. The normalization of these CS's is
given by
\begin{align}
\left\langle \left.  \Phi\left(  \mathbf{z}\right)  \right\vert \Phi\left(
\mathbf{z}\right)  \right\rangle  &  =\sum_{1\leq j\leq\left(
%TCIMACRO{\QATOP{r}{N}}%
%BeginExpansion
\genfrac{}{}{0pt}{}{r}{N}%
%EndExpansion
\right)  }\left\vert z_{j}\right\vert ^{2}=\left\vert z_{1}\right\vert
^{2}+\sum_{2\leq j\leq\left(
%TCIMACRO{\QATOP{r}{N}}%
%BeginExpansion
\genfrac{}{}{0pt}{}{r}{N}%
%EndExpansion
\right)  }\left\vert z_{j}\right\vert ^{2}\\
\left\langle \left.  \Phi\left(  \mathbf{\eta}\right)  \right\vert \Phi\left(
\mathbf{\eta}\right)  \right\rangle  &  =e^{2\operatorname{Re}\eta_{1}%
}\left\{  1+\sum_{2\leq j\leq\left(
%TCIMACRO{\QATOP{r}{N}}%
%BeginExpansion
\genfrac{}{}{0pt}{}{r}{N}%
%EndExpansion
\right)  }\left\vert \eta_{j}\right\vert ^{2}\right\}
\end{align}
Hence if one sets $\eta_{1}$ as the function
\begin{equation}
\eta_{1}=\eta_{1}\left(  \eta_{2},\cdots,\eta_{\left(
%TCIMACRO{\QATOP{r}{N}}%
%BeginExpansion
\genfrac{}{}{0pt}{}{r}{N}%
%EndExpansion
\right)  }\right)  =\frac{1}{2}\ln\left(  1-\sum_{2\leq j\leq\left(
%TCIMACRO{\QATOP{r}{N}}%
%BeginExpansion
\genfrac{}{}{0pt}{}{r}{N}%
%EndExpansion
\right)  }e^{2\operatorname{Re}\eta_{j}}\right)  \label{n1}%
\end{equation}
and $z_{1}$ as the function%
\begin{equation}
z_{1}=z_{1}\left(  z_{2},\cdots,z_{\left(
%TCIMACRO{\QATOP{r}{N}}%
%BeginExpansion
\genfrac{}{}{0pt}{}{r}{N}%
%EndExpansion
\right)  }\right)  =\left(  1-\sum_{2\leq j\leq\left(
%TCIMACRO{\QATOP{r}{N}}%
%BeginExpansion
\genfrac{}{}{0pt}{}{r}{N}%
%EndExpansion
\right)  }\left\vert z_{j}\right\vert ^{2}\right)  ^{\frac{1}{2}} \label{n2}%
\end{equation}
the family of CS's $\left\{  \left\vert \Phi\left(  \mathbf{z}\right)
\right\rangle \right\}  \equiv\left\{  \left\vert \Phi\left(  \mathbf{\eta
}\right)  \right\rangle \right\}  $ are normalized if the reference state
$\left\vert \Phi_{1}\right\rangle $ is normalized i.e. $\left\langle \left.
\Phi_{1}\right\vert \Phi_{1}\right\rangle =1.$ If one uses the CS parameters
$\left\{  \eta_{2},\cdots,\eta_{\left(
%TCIMACRO{\QATOP{r}{N}}%
%BeginExpansion
\genfrac{}{}{0pt}{}{r}{N}%
%EndExpansion
\right)  }\right\}  $ or $\left\{  z_{2},\cdots,z_{\left(
%TCIMACRO{\QATOP{r}{N}}%
%BeginExpansion
\genfrac{}{}{0pt}{}{r}{N}%
%EndExpansion
\right)  }\right\}  $ the CS manifold is the unit sphere in $\mathcal{H}^{N}$,
while if one uses $\left\{  \mathbf{\eta}\right\}  $ or $\left\{
\mathbf{z}\right\}  $ then it is all of $\mathcal{H}^{N}$ and one explicitly
considers the normalization in all expressions.

\section{Discussion\label{discussion}}

\appendix

\section{Conjecture\label{conjecture}}

\section{AGP States \label{agp}}

\bibliographystyle{apsrev}
\bibliography{agpopt,bibrefs,bw,Coleman,lowdin-per-olov,novel,Ohrn,Osvaldo,Yngve}

\end{document}